\documentstyle[%

righttag%

,11pt%

,amssymb%

,russian%

,izvru%

]{amsart}



\setcounter{page}{85} 

\pagestyle{plain}

\begin{document}



\begin{center}

{\Large \bf Информация о научной конференции\\ ``Геометрическая 

теория функций и краевые задачи''}

\end{center}

\vskip 1cm



Казанским государственным университетом на базе НИИ математики и

механики им. Н.Г.\,Че\-ботарева 18--22 марта 2002 г. в г.\,Казани

была проведена научная конференция ``Геометрическая 

теория функций и краевые задачи'', посвященная 200-летию Казанского

университета. 



Данный тематический выпуск содержит новые математические 

результаты, доложенные на этой конференции.

Публикация работ

будет продолжена в \No\,2 журнала за 2003 год.



На конференции в качестве лекторов выступили крупные специалисты 

по теории функций. Среди них член-корреспондент РАН П.Л.\,Ульянов

(Москва), профессора Л.А.\,Бекларян, И.Х.\,Сабитов (Москва), 

С.К.\,Водопьянов (Новосибирск), Д.В.\,Прохоров, Г.В.\,Хромова (Саратов), 

К.М.\,Расулов (Смоленск), Э.И.\,Зверович (Минск), В.В.\,Старков

(Петрозаводск), Р.С.\,Юл\-му\-ха\-ме\-тов (Уфа), Э.Г.\,Кирьяцкий 

(Вильнюс). Всего сделано 45 докладов. В конференции приняли участие 

с докладами и лекциями: доктора наук, профессора --- 34, кандидаты 

наук, доценты --- 34, аспиранты, молодые ученые, студенты --- 19. 

По городам: Москва --- 5, Саратов --- 5, Минск --- 4, Йошкар-Ола --- 4, 

Ростов-на-Дону --- 2, Вильнюс --- 2, Уфа --- 4, Сургут, Краснодар, 

Ульяновск, Нижний Новгород, Караганда, Волгоград, Петрозаводск, 

Пермь, Брянск, Новосибирск, Кустанай, Самара, Смоленск --- по 1, 

Казань --- 40.



Конференция проведена при участии Казанского математического 

общества и Математического центра им. Н.И.\,Лобачевского.

Финансовую поддержку оказали Российский фонд фундаментальных 

исследований и Академия наук Республики Татарстан. 

Расширенные тезисы докладов участников конференции составили 

13-й том ``Трудов Математического центра им.~Н.И.\,Лобачевского''

(Казань: Изд-во Казанск. матем. о-ва, 2002. -- 182 с.)

Материалы, 

представленные приглашенными докладчиками, составили 

содержание 14-го тома ``Трудов Математического центра 

им.~Н.И.Лобачевского''

(Казань: Изд-во Казанск. матем. о-ва, 2002. -- 368 с.) и 

настоящего тематического выпуска журнала ``Известия вузов. 

Математика''.



Участники конференции тепло и сердечно поздравили профессора

Казанского университета Леонида Александровича Аксентьева

с 70-летием и 40-летием руководимого им в г.\,Казани

Городского семинара по геометрической теории функций.



Редколлегия и редакция журнала ``Известия вузов. Математика'' также

сердечно поздравляют Леонида Александровича Аксентьева,

заместителя главного редактора журнала, с 70-летием

и желают ему успехов в научной работе, здоровья на долгие лета,

счастья!



\end{document}













